\section{Introduction}
\label{sec:intro}

Many materials comprise networks of long, overlapping strands, including
carbon-nanotube (CNT) sheets, collagen in connective tissue, paper and nonwoven
textiles. Their stiffness, strength and conductivity depend on network
geometry, constituent behaviour and inter-strand interactions. For example,
load transfer between nanotubes can limit an assembly's stiffness and strength
despite the high modulus of each tube
\citep{devolder2013cnt,coleman2006small,volkov2010structural,%
ostanin2013distinct,ostanin2019size,wittmaack2018mesoscopic}.
Predictive models therefore require geometric inputs such as strand length,
orientation, curvature and connectivity. Bulk image measurements such as
coverage and orientation do not establish which visible segments belong to
the same strand \citep{dong2022detecting,imtiaz2023facile,beigzadeh2024utilizing}.
Recovering that identity is necessary to measure individual strands end to end.

This is an instance-segmentation problem in which elongated, open curves may
share pixels at crossings. Crossings merge distinct strands into one mask
component, whereas signal gaps split one strand into fragments whose tips can
be indistinguishable from genuine ends. Reconstruction must therefore decide
which segments continue through each junction and which reconnect across a
gap. Here, we address this task from a binary strand-axis mask, whose foreground
pixels mark strand centrelines.

Existing methods recover strand continuity using local pairing rules, learned
image features or joint optimisation. However, greedy pairing can commit to
locally plausible but incompatible continuations, while alternatives may
require annotated training data, restrict the number of segments considered
at a junction, or depend on anatomical seeds
\citep{zhang2017sifne,liu2019intersection,playout2026dnai,de2016graph}.
The remaining need is to resolve competing continuations at complex junctions
from a mask alone, while allowing genuine ends, reconnecting gaps and retaining
shared crossing pixels. Section~\ref{sec:related} reviews these approaches and
their relation to the present work.

PLECTA addresses this need by jointly selecting pairs of arms, the strand
segments entering a junction, while allowing arms to remain unpaired. It
solves an exact maximum-weight \emph{partial}-matching problem at each junction,
with an explicit penalty for unmatched arms and no restriction on junction
degree. It combines this decision with geometric checks for reconnecting gaps
and updates to direction and curvature estimates as longer strands are
assembled. The output allows crossing pixels to belong to multiple strands.
These ingredients are individually established; their combination is the
contribution of this work. Reconstruction is deterministic for fixed inputs
and parameters and assumes directional continuity, targeting semiflexible
networks whose strands maintain their direction locally through crossings.

We evaluate PLECTA on synthetic scenes and annotated scanning electron
microscopy (SEM) images of CNT sheets, using paired clean and degraded masks
to isolate the effect of mask quality. Across the evaluated datasets, PLECTA
reconstructed strand identity more accurately than the compared methods,
although performance remained sensitive to network density and mask quality.
For SEM images, an upstream segmenter supplies the strand-axis masks
\citep{ruehle2021workflow}. A separate stage uses the corresponding greyscale
image to estimate which strand lies on top at each crossing; this is a proof
of concept for depth ordering, not a validated three-dimensional
reconstruction. Section~\ref{sec:methods} describes the workflow and evaluation,
Section~\ref{sec:results} reports the results, and
Section~\ref{sec:conclusions} discusses conclusions and limitations.
